\documentclass{ximera}

\input{../../../preamble.tex}


\definecolor{fvdnavy}{HTML}{071D43}
\definecolor{fvdcyan}{HTML}{20BFC6}
\definecolor{fvdcyanlight}{HTML}{EFFCFB}
\definecolor{fvdmagenta}{HTML}{F10D69}
\definecolor{fvdmagentalight}{HTML}{FFF1F7}
\definecolor{fvdtext}{HTML}{163044}





%=================================================
% Pedagogische omgevingen
% PDF: tcolorbox
% HTML: learning-box
%=================================================

\newenvironment{voorkennis}
{%
    \ifdefined\HCode
        \HCode{%
<section class="learning-box learning-box--prior">
<div class="learning-box__header">
<span class="learning-box__icon" aria-hidden="true"></span>
<span class="learning-box__title">Voorkennis</span>
</div>
<div class="learning-box__content">
        }%
        \begin{itemize}
    \else
       \begin{tcolorbox}[
    enhanced,
    breakable,
    colback=fvdcyanlight,
    colframe=fvdcyan,
    coltext=fvdtext,
    boxrule=0.5pt,
    leftrule=4pt,
    arc=4mm,
    outer arc=4mm,
    left=7mm,
    right=6mm,
    top=5mm,
    bottom=5mm,
    before skip=6mm,
    after skip=6mm,
    title=Voorkennis,
    coltitle=fvdcyan,
    fonttitle=\bfseries\Large,
    attach title to upper={\par\smallskip},
    boxed title style={
        empty,
        left=0mm,
        right=0mm,
        top=0mm,
        bottom=0mm
    },
    overlay={
        \draw[
            fvdcyan,
            line width=0.5pt
        ]
        ([xshift=42mm,yshift=-13mm]frame.north west)
        --
        ([xshift=-7mm,yshift=-13mm]frame.north east);
    }
]
        \begin{itemize}
    \fi
}
{%
    \end{itemize}
    \ifdefined\HCode
        \HCode{%
</div>
</section>
        }%
    \else
        \end{tcolorbox}
    \fi
}


\newenvironment{leerdoelen}
{%
    \ifdefined\HCode
        \HCode{%
<section class="learning-box learning-box--goals">
<div class="learning-box__header">
<span class="learning-box__icon" aria-hidden="true"></span>
<span class="learning-box__title">Leerdoelen</span>
</div>
<div class="learning-box__content">
        }%
        \begin{itemize}
    \else
\begin{tcolorbox}[
    enhanced,
    breakable,
    colback=fvdmagentalight,
    colframe=fvdmagenta,
    coltext=fvdtext,
    boxrule=0.5pt,
    leftrule=4pt,
    arc=4mm,
    outer arc=4mm,
    left=7mm,
    right=6mm,
    top=5mm,
    bottom=5mm,
    before skip=6mm,
    after skip=6mm,
    title=Leerdoelen,
    coltitle=fvdmagenta,
    fonttitle=\bfseries\Large,
    attach title to upper={\par\smallskip},
    boxed title style={
        empty,
        left=0mm,
        right=0mm,
        top=0mm,
        bottom=0mm
    },
    overlay={
        \draw[
            fvdmagenta,
            line width=0.5pt
        ]
        ([xshift=42mm,yshift=-13mm]frame.north west)
        --
        ([xshift=-7mm,yshift=-13mm]frame.north east);
    }
]
        \begin{itemize}
    \fi
}
{%
    \end{itemize}
    \ifdefined\HCode
        \HCode{%
</div>
</section>
        }%
    \else
        \end{tcolorbox}
    \fi
}


\addPrintStyle{../../..}

\title{Oefeningen — Stelsels aan het werk}
\author{Ingmar Herreman}

\begin{document}

% FVD_CHAPTER_DATA_START
% Automatisch gegenereerd uit index.tex.
% Niet handmatig aanpassen.
\fvdchapterdata{1}{Van algebra naar ruimte}{2}
% FVD_CHAPTER_DATA_END





\maketitle

\FVDbeginExercises

\begin{opmerking}
Gebruik bij de oefeningen ICT voor de matrixreductie, tenzij uitdrukkelijk
anders vermeld.

Noteer bij een volledig modelleerprobleem steeds:
\begin{enumerate}
  \item de betekenis van de onbekenden;
  \item het stelsel;
  \item de uitgebreide matrix;
  \item de relevante ICT-uitvoer;
  \item de interpretatie van de oplossing in de context;
  \item een korte controle.
\end{enumerate}
\end{opmerking}


% ============================================================
% BASIS
% ============================================================

\FVDsection{Basis — van context naar oplossing}

\begin{oefening}[Bloemenprijzen]{\oefB\ESS}

In de lente worden \(24\) geraniums, \(12\) hanggeraniums en \(18\)
lobelia's gekocht voor \(47{,}40\) euro.

Voor een grote bloembak worden \(6\) geraniums, \(3\) hanggeraniums en
\(2\) lobelia's gebruikt. Deze bloemen kosten samen \(10{,}60\) euro.

Voor een pergola worden \(5\) geraniums, \(7\) hanggeraniums en \(6\)
lobelia's gebruikt. Deze kosten samen \(16{,}40\) euro.

Bepaal met een lineair stelsel en ICT de prijs per stuk van elke bloemsoort.

\begin{oplossing}
Neem \(x,y,z\) als de prijs in euro van respectievelijk een geranium,
hanggeranium en lobelia.

\[
\begin{cases}
24x+12y+18z=47{,}40,\\
6x+3y+2z=10{,}60,\\
5x+7y+6z=16{,}40.
\end{cases}
\]

ICT geeft:
\[
x=1,\qquad y=1{,}20,\qquad z=0{,}50.
\]

Dus een geranium kost \(1{,}00\) euro, een hanggeranium \(1{,}20\) euro
en een lobelia \(0{,}50\) euro.
\end{oplossing}
\end{oefening}


\begin{oefening}[Voedingsmengsel]{\oefB\ESS}

Een bioloog wil een voedselmengsel maken dat \(23\) g proteïne,
\(6{,}2\) g vet en \(16\) g vocht bevat.

\begin{itemize}
  \item A bevat \(20\%\) proteïne, \(2\%\) vet en \(15\%\) vocht.
  \item B bevat \(10\%\) proteïne, \(6\%\) vet en \(10\%\) vocht.
  \item C bevat \(15\%\) proteïne, \(5\%\) vet en \(5\%\) vocht.
\end{itemize}

Hoeveel gram van elk mengsel is nodig?

\begin{oplossing}
Neem \(x,y,z\) als de massa's in gram.

\[
\begin{cases}
0{,}20x+0{,}10y+0{,}15z=23,\\
0{,}02x+0{,}06y+0{,}05z=6{,}2,\\
0{,}15x+0{,}10y+0{,}05z=16.
\end{cases}
\]

ICT geeft
\[
(x,y,z)=(60,50,40).
\]

Er is dus \(60\) g A, \(50\) g B en \(40\) g C nodig.
\end{oplossing}
\end{oefening}


\begin{oefening}[Een parabool door drie punten]{\oefB\ESS}

Een parabool heeft voorschrift
\[
y=ax^2+bx+c
\]
en gaat door de punten
\[
(1,1),\qquad (-1,5),\qquad (2,2).
\]

Stel een lineair stelsel op voor \(a,b,c\), los het met ICT op en bepaal
het voorschrift.

\begin{oplossing}
Invullen van de drie punten geeft
\[
\begin{cases}
a+b+c=1,\\
a-b+c=5,\\
4a+2b+c=2.
\end{cases}
\]

ICT geeft
\[
a=1,\qquad b=-2,\qquad c=2.
\]

Dus
\[
\boxed{y=x^2-2x+2}.
\]
\end{oplossing}
\end{oefening}


\begin{oefening}[Tuinmest mengen]{\oefB\ESS}

Een groothandel beschikt over drie meststoffen:

\begin{itemize}
  \item A bevat \(5\%\) stikstof;
  \item B bevat \(10\%\) stikstof;
  \item C bevat \(20\%\) stikstof.
\end{itemize}

Er moet \(1000\) kg mest met \(15\%\) stikstof worden gemaakt.
Van soort C wil men tweemaal zoveel gebruiken als van soort A.

Bepaal hoeveel kilogram van elke soort nodig is.

\begin{oplossing}
Neem \(x,y,z\) als de massa's A, B en C.

\[
\begin{cases}
x+y+z=1000,\\
0{,}05x+0{,}10y+0{,}20z=150,\\
z=2x.
\end{cases}
\]

ICT geeft
\[
x=\frac{1000}{3},\qquad y=0,\qquad z=\frac{2000}{3}.
\]

Dus ongeveer
\[
333{,}333\text{ kg A},\qquad
0\text{ kg B},\qquad
666{,}667\text{ kg C}.
\]

Controle:
\[
0{,}05\cdot\frac{1000}{3}
+
0{,}20\cdot\frac{2000}{3}
=150.
\]
\end{oplossing}
\end{oefening}


\begin{oefening}[Lees eerst, reken daarna]{\oefB\ESS}

Een laboratorium mengt drie oplossingen A, B en C. De massa's zijn
\(x,y,z\) gram. Uit de context werden de volgende voorwaarden afgeleid:
\[
\begin{cases}
x+y+z=200,\\
0{,}10x+0{,}20y+0{,}30z=42,\\
x-z=20.
\end{cases}
\]

\begin{question}
Schrijf de uitgebreide matrix op.
\begin{oplossing}
\[
\left[
\begin{array}{ccc|c}
1&1&1&200\\
0{,}10&0{,}20&0{,}30&42\\
1&0&-1&20
\end{array}
\right].
\]
\end{oplossing}
\end{question}

\begin{question}
Reduceer met ICT en interpreteer de oplossing.
\begin{oplossing}
ICT geeft
\[
x=80,\qquad y=60,\qquad z=60.
\]
Er is \(80\) g A, \(60\) g B en \(60\) g C nodig.
\end{oplossing}
\end{question}
\end{oefening}


% ============================================================
% VERDIEPING
% ============================================================

\FVDsection{Verdieping — modelleren en interpreteren}

\begin{oefening}[Zestig ledlampen]{\oefU\ESS}

Een ruimte moet worden verlicht met precies \(60\) ledlampen en een totale
lichtstroom van \(40\,000\) lumen. Er zijn lampen van \(900\), \(500\) en
\(400\) lumen.

\begin{question}
Kies onbekenden en stel het stelsel op.
\begin{oplossing}
Met \(x,y,z\) als de aantallen lampen:
\[
\begin{cases}
x+y+z=60,\\
900x+500y+400z=40\,000.
\end{cases}
\]
\end{oplossing}
\end{question}

\begin{question}
Gebruik ICT. Leg met rang uit waarom het stelsel niet één unieke
algebraïsche oplossing heeft.
\begin{oplossing}
Er zijn drie onbekenden en de twee vergelijkingen zijn onafhankelijk:
\[
\operatorname{rang}(A)=\operatorname{rang}(A_b)=2<3.
\]
Er is dus één vrije parameter en algebraïsch zijn er oneindig veel oplossingen.
\end{oplossing}
\end{question}

\begin{question}
Bepaal alle praktisch mogelijke combinaties.
\begin{oplossing}
De niet-negatieve gehele oplossingen zijn:
\[
\begin{array}{c|c|c}
x&y&z\\ \hline
25&35&0\\
26&30&4\\
27&25&8\\
28&20&12\\
29&15&16\\
30&10&20\\
31&5&24\\
32&0&28
\end{array}
\]
De context reduceert de oneindige reële oplossingsverzameling dus tot acht
praktisch mogelijke keuzes.
\end{oplossing}
\end{question}
\end{oefening}


\begin{oefening}[Een onmogelijk productieplan]{\oefU\ESS}

Een atelier produceert drie onderdelen \(A,B,C\). Voor een dagplanning
worden de volgende voorwaarden opgesteld:
\[
\begin{cases}
x+y+z=100,\\
2x+y+3z=180,\\
4x+2y+6z=365.
\end{cases}
\]

Gebruik ICT en analyseer de oplosbaarheid. Wat betekent je besluit voor
het productieplan?

\begin{oplossing}
De linkerkant van de derde vergelijking is tweemaal de linkerkant van
de tweede, maar
\[
365\neq2\cdot180.
\]
Na reductie ontstaat een strijdige rij. Dus
\[
\operatorname{rang}(A)<\operatorname{rang}(A_b).
\]
Er bestaat geen productieplan dat aan alle drie de voorwaarden tegelijk voldoet.
Minstens één opgelegde voorwaarde of één gegeven moet worden herzien.
\end{oplossing}
\end{oefening}


\begin{oefening}[Is de oplossing bruikbaar?]{\oefU\ESS}

Een stelsel voor de aantallen van drie types sensoren levert met ICT
\[
(x,y,z)=(12{,}5,8,19{,}5).
\]

Een leerling besluit: ``Het probleem is opgelost.''

Beoordeel die uitspraak.

\begin{oplossing}
Algebraïsch kan de oplossing correct zijn, maar aantallen sensoren moeten
gehele, niet-negatieve getallen zijn. De waarden \(12{,}5\) en \(19{,}5\)
zijn dus niet uitvoerbaar.

De juiste conclusie is dat het lineaire model een unieke reële oplossing
heeft, maar geen bruikbare oplossing binnen de context van aantallen sensoren.
\end{oplossing}
\end{oefening}


\begin{oefening}[Welke vergelijking hoort bij welke voorwaarde?]{\oefU}

Een bedrijf maakt drie producten \(P,Q,R\). Voor één product zijn
respectievelijk de volgende hoeveelheden grondstof en machinetijd nodig:

\[
\begin{array}{c|ccc}
& P&Q&R\\ \hline
\text{grondstof (kg)}&2&3&1\\
\text{machinetijd (min)}&4&2&5
\end{array}
\]

Er is \(200\) kg grondstof en \(360\) minuten machinetijd beschikbaar.
Men wil in totaal \(100\) producten maken.

\begin{question}
Stel het stelsel op.
\begin{oplossing}
Met \(x,y,z\) als de aantallen \(P,Q,R\):
\[
\begin{cases}
x+y+z=100,\\
2x+3y+z=200,\\
4x+2y+5z=360.
\end{cases}
\]
\end{oplossing}
\end{question}

\begin{question}
Los met ICT op en controleer of het plan praktisch mogelijk is.
\begin{oplossing}
ICT geeft
\[
x=20,\qquad y=40,\qquad z=40.
\]
Alle aantallen zijn niet-negatieve gehele getallen. Controle:
\[
20+40+40=100,
\]
\[
2(20)+3(40)+40=200,
\]
Ook de grondstof- en machinetijdvoorwaarden zijn voldaan. Het productieplan is dus praktisch mogelijk.
\end{oplossing}
\end{question}
\end{oefening}


% ============================================================
% CHEMIE
% ============================================================

\FVDsection{STEM — reactievergelijkingen balanceren}

\begin{oefening}[Ammoniak oxideert]{\oefB\ESS}

Balanceer met een homogeen stelsel en ICT:
\[
\ldots\,NH_3+\ldots\,O_2
\longrightarrow
\ldots\,NO+\ldots\,H_2O.
\]

Noteer ook de rang, het aantal vrije parameters en verklaar waarom je
uiteindelijk gehele coëfficiënten kiest.

\begin{oplossing}
Neem
\[
aNH_3+bO_2\longrightarrow cNO+dH_2O.
\]
Atoombehoud geeft
\[
\begin{cases}
a-c=0,\\
3a-2d=0,\\
2b-c-d=0.
\end{cases}
\]
Er zijn vier onbekenden en drie onafhankelijke vergelijkingen:
\[
\operatorname{rang}(A)=3.
\]
Er is dus één vrije parameter. De kleinste positieve gehele verhouding is
\[
(a,b,c,d)=(4,5,4,6).
\]
Dus
\[
\boxed{4NH_3+5O_2\longrightarrow4NO+6H_2O}.
\]
\end{oplossing}
\end{oefening}


\begin{oefening}[Neerslagreactie]{\oefB\ESS}

Balanceer met een stelsel:
\[
\ldots\,AgNO_3+\ldots\,HCl
\longrightarrow
\ldots\,AgCl+\ldots\,HNO_3.
\]

\begin{oplossing}
Met coëfficiënten \(a,b,c,d\) volgt uit atoombehoud een homogeen stelsel.
De kleinste positieve gehele oplossing is
\[
(a,b,c,d)=(1,1,1,1).
\]
Dus
\[
\boxed{AgNO_3+HCl\longrightarrow AgCl+HNO_3}.
\]
\end{oplossing}
\end{oefening}


\begin{oefening}[IJzer(III)hydroxide en zwavelzuur]{\oefU\ESS}

Balanceer volledig met een homogeen stelsel en ICT:
\[
\ldots\,Fe(OH)_3+\ldots\,H_2SO_4
\longrightarrow
\ldots\,Fe_2(SO_4)_3+\ldots\,H_2O.
\]

\begin{oplossing}
Neem
\[
aFe(OH)_3+bH_2SO_4\longrightarrow cFe_2(SO_4)_3+dH_2O.
\]
Atoombehoud levert onder andere
\[
a=2c,\qquad b=3c,
\]
en uit waterstof
\[
3a+2b=2d.
\]
Kies \(c=1\). Dan
\[
a=2,\qquad b=3,\qquad d=6.
\]
Dus
\[
\boxed{2Fe(OH)_3+3H_2SO_4\longrightarrow Fe_2(SO_4)_3+6H_2O}.
\]
\end{oplossing}
\end{oefening}


\begin{oefening}[Verbranding van propaan]{\oefB\ESS}

Balanceer:
\[
\ldots\,C_3H_8+\ldots\,O_2
\longrightarrow
\ldots\,CO_2+\ldots\,H_2O.
\]

Gebruik expliciet de stappen
\[
\text{atoombehoud}\to\text{stelsel}\to\text{matrix}\to\text{ICT}
\to\text{gehele verhouding}.
\]

\begin{oplossing}
Neem coëfficiënten \(a,b,c,d\). Dan
\[
\begin{cases}
3a-c=0,\\
8a-2d=0,\\
2b-2c-d=0.
\end{cases}
\]
De kleinste positieve gehele verhouding is
\[
(a,b,c,d)=(1,5,3,4).
\]
Dus
\[
\boxed{C_3H_8+5O_2\longrightarrow3CO_2+4H_2O}.
\]
\end{oplossing}
\end{oefening}


\begin{oefening}[Roosten van pyriet]{\oefU}
Pyriet \(FeS_2\) kan reageren met zuurstof:
\[
\ldots\,FeS_2+\ldots\,O_2
\longrightarrow
\ldots\,Fe_2O_3+\ldots\,SO_2.
\]

Balanceer met een homogeen stelsel.

\begin{oplossing}
Met coëfficiënten \(a,b,c,d\):
\[
\begin{cases}
a-2c=0,\\
2a-d=0,\\
2b-3c-2d=0.
\end{cases}
\]
De kleinste positieve gehele verhouding is
\[
(a,b,c,d)=(4,11,2,8).
\]
Dus
\[
\boxed{4FeS_2+11O_2\longrightarrow2Fe_2O_3+8SO_2}.
\]
\end{oplossing}
\end{oefening}


\begin{oefening}[Waarom is \(2,4,2,4\) niet het eindantwoord?]{\oefU\ESS}

Een leerling balanceert methaanverbranding en vindt
\[
2CH_4+4O_2\longrightarrow2CO_2+4H_2O.
\]

De vergelijking is qua atoomaantallen correct. Toch wordt gewoonlijk
\[
CH_4+2O_2\longrightarrow CO_2+2H_2O
\]
genoteerd.

Leg dit uit vanuit de oplossingsverzameling van het homogene stelsel.

\begin{oplossing}
Alle veelvouden van een niet-triviale oplossing zijn opnieuw oplossingen
van het homogene stelsel. De vector
\[
(2,4,2,4)
\]
is tweemaal
\[
(1,2,1,2).
\]
Beide geven dezelfde verhouding. Conventioneel kiest men de kleinste
positieve gehele coëfficiënten.
\end{oplossing}
\end{oefening}


% ============================================================
% GEVORDERD
% ============================================================

\FVDsection{Gevorderd — kritisch modelleren}

\begin{oefening}[Ontbrekende informatie]{\oefG}

Een onderzoeker wil drie oplossingen A, B en C mengen. Hij weet alleen:
\[
x+y+z=500
\]
en
\[
0{,}10x+0{,}20y+0{,}30z=100.
\]

\begin{question}
Analyseer met rang hoeveel vrijheid het model nog bevat.
\begin{oplossing}
Er zijn drie onbekenden en twee onafhankelijke vergelijkingen:
\[
\operatorname{rang}(A)=\operatorname{rang}(A_b)=2<3.
\]
Er blijft één vrije parameter over. De gegevens bepalen dus niet één uniek mengsel.
\end{oplossing}
\end{question}

\begin{question}
Geef een voorbeeld van een extra lineaire voorwaarde die tot een unieke
oplossing kan leiden.
\begin{oplossing}
Bijvoorbeeld
\[
x=z
\]
of een opgelegde hoeveelheid van één mengsel. De extra voorwaarde moet
onafhankelijk zijn van de twee bestaande voorwaarden.
\end{oplossing}
\end{question}
\end{oefening}


\begin{oefening}[Ontwerp zelf een stelsel]{\oefG}
  
Bedenk zelf een realistische STEM-context met drie onbekenden en minstens
drie lineaire voorwaarden.

\begin{enumerate}
  \item Beschrijf de context en de betekenis van de onbekenden.
  \item Stel het stelsel op.
  \item Voorspel vóór de ICT-berekening of je één, geen of mogelijk
        oneindig veel oplossingen verwacht.
  \item Reduceer de uitgebreide matrix met ICT.
  \item Bepaal de relevante rangen.
  \item Interpreteer en controleer de oplossing.
  \item Geef aan welke aannames of beperkingen jouw model heeft.
\end{enumerate}

\begin{oplossing}
Eigen antwoord. Een volledig antwoord bevat een consistente context,
een correct lineair stelsel, een correcte ICT-analyse en een expliciete
terugkoppeling naar de context.
\end{oplossing}
\end{oefening}


% ============================================================
% EINDCHECK
% ============================================================

\FVDsection{Eindcheck}

\begin{oefening}[Van tekst tot besluit]{\oefU\ESS}

Een labo wil \(100\) ml van een oplossing maken uit drie voorraadoplossingen
A, B en C. Hun concentraties zijn respectievelijk \(10\%\), \(20\%\) en
\(40\%\). Het eindmengsel moet \(25\%\) bevatten. Bovendien moet de
hoeveelheid A gelijk zijn aan de hoeveelheid C.

Los het probleem volledig op met een stelsel en ICT.

\begin{oplossing}
Neem \(x,y,z\) als de volumes in ml.

\[
\begin{cases}
x+y+z=100,\\
0{,}10x+0{,}20y+0{,}40z=25,\\
x-z=0.
\end{cases}
\]

Omdat \(z=x\), volgt
\[
2x+y=100
\]
en
\[
0{,}50x+0{,}20y=25.
\]
ICT geeft
\[
x=z=\frac{50}{3},\qquad y=\frac{200}{3}.
\]

Dus ongeveer
\[
16{,}67\text{ ml A},\qquad
66{,}67\text{ ml B},\qquad
16{,}67\text{ ml C}.
\]

Alle volumes zijn positief en samen \(100\) ml. De concentratiecontrole
geeft \(25\) ml werkzame stof per \(100\) ml mengsel.
\end{oplossing}
\end{oefening}



\end{document}
